\section{Introduction}

The foundational role of group theory in Quantum Field Theory (QFT) was
established by Wigner \cite{Wigner1939}, who showed that physical states are
defined as unitary irreducible representations of the inhomogeneous Lorentz
group (Poincar\'{e} group). This classification, based on the ``little group''
method, reveals that fundamental wave equations emerge as specific realizations
of the underlying symmetry structure: scalar fields (spin 0) are governed by the
Klein-Gordon equation, while vector fields (spin 1) correspond to the Maxwell
and Proca equations \cite{ryderQuantumFieldTheory1996}. Within this framework, the Dirac equation
occupies a role for describing fermionic fields; while the modern
group-theoretical approach derives it directly from the transformation
properties of spinors under the Lorentz group, its original formulation was
motivated by the necessity to resolve ``duplexity'' phenomena in atomic spectra,
leading Dirac to construct a wave equation linear in time and momentum that
satisfied the requirements of both relativity and quantum mechanics
\cite{diracQuantumTheoryElectron1928}.

In recent developments, Watson and Musielak \cite{watsonChiralSymmetryDirac2020,
watsonChiralDiracEquation2021, Watson2022} revisited the
group-theoretical foundations of the field description. Instead of relying on
the standard factorization of the Klein-Gordon equation, they derived the field
equations directly from the irreducible representations of the Poincar\'{e}
group extended by parity. This approach revealed an intrinsic freedom in the
choice of the chiral basis, parametrized by a new physical degree of freedom
identified as the chiral angle. By identifying this degree of freedom, they
obtained a generalized form of the Dirac equation, hereafter referred to as the
Chiral Dirac Equation (CDE), which incorporates a pseudoscalar mass term
alongside the standard scalar mass. The authors suggested that this additional
chiral degree of freedom has physical implications; specifically, the interplay
between the scalar and pseudoscalar mass components offers a mechanism to
explain the suppression of neutrino masses, proposing the new massive particle
as a candidate for Dark Matter.

The first critical issue concerns the consistency of the quantum mechanical
description derived from the CDE. Standard quantum mechanics postulates that
physical observables must be represented by self-adjoint operators to guarantee
real energy eigenvalues and unitary time evolution (probability conservation).
While the Dirac equation is consistent with these principles, the introduction
of a chiral angle and a pseudoscalar mass term modifies the structure of the
evolution operator. It is not immediately evident that the Hamiltonian
associated with this generalized equation retains hermiticity under the new
chiral basis transformations. Therefore, we aim to explicitly derive the
Hamiltonian from the CDE and test its hermiticity to ensure the theory's
physical viability.

Secondly, a significant algebraic inconsistency arises when applying the
framework to neutrino physics under strict chirality constraints \cite{watsonChiralSymmetryDirac2020}.
Preliminary analyses indicate that for a pure left-handed field, defined by the chiral
projection $\nu_L = \frac{1}{2}\left(1-\gamma^5\right)\nu$, the scalar bilinear
$\overline{\nu}\nu$ vanishes identically. Consequently, the Yukawa Lagrangian
$\mathcal{L}_Y$ derived in the original work becomes trivial
($\mathcal{L}_Y = 0$), thereby nullifying the proposed mechanism for generating
massive neutrinos in the Dirac context. To address this, we propose to explore
the Majorana mass term within the context of a more general see-saw mechanism,
functioning as a chiral see-saw.

Finally, while the original derivation focused on bispinors, effectively
restricting the analysis to spin-1/2 particles
\cite{watsonChiralSymmetryDirac2020}, Watson et al. have also explored
vector bosons \cite{Watson2022}. However, the method employed---based on the irreducible
representations of the Poincar\'{e} group---is fundamentally geometric and
should be applicable to other spin sectors. To this end, we propose a
review of the feasibility of extending this chiral formalism to
higher-spin fields. Specifically, we examine whether similar chiral degrees of
freedom can be consistently introduced into the Proca equation for massive
vector bosons (spin-1) or the Rarita-Schwinger fields (spin-3/2).
